% ==========================================
% HAUPTKAPITEL: Software und Elektronik
% ==========================================
\chapter{Software und Elektronik}
% Falls euer Dokument auf \section-Ebene beginnt, ändere \chapter einfach zu \section
% und schiebe die folgenden Gliederungspunkte entsprechend noch eine Ebene tiefer.

% ------------------------------------------
% 1. Requirements Engineering
% ------------------------------------------
\section{Requirements Engineering}
\begin{table}[h]
	\centering
	\caption{Auszug der Hardware- und Softwareanforderungen (Software \& Elektronik)}
	\label{tab:requirements}
	% tabularx nutzt die volle Textbreite (\textwidth). 
	% Die Spalten: l = linksbündig, X = automatische Breite mit Zeilenumbruch, c = zentriert
	\begin{tabularx}{\textwidth}{l X c}
		\toprule
		\textbf{ID} & \textbf{Anforderung (Beschreibung)} & \textbf{Typ} \\
		\midrule
		ID-0008 & Die Software muss die zwei Temperatursensoren über ADC-Kanäle auslesen und die Werte in Echtzeit verarbeiten. & Funktional \\
		\addlinespace
		ID-0010 & Die Software muss Eingaben vom Drehregler über einen ADC-Kanal oder EXTI-Interrupt erfassen und interpretieren. & Funktional \\
		\addlinespace
		ID-0011 & Die Software muss das LCD-Display über I²C ansteuern und relevante Informationen wie die aktuellen Temperaturen anzeigen. & Funktional \\
		\addlinespace
		ID-0013 & Die Software muss PWM-Signale generieren, um die Pumpe und den Motor stufenlos zu steuern (anstelle eines Relais). & Funktional \\
		\addlinespace
		ID-0014 & Die Software muss DMA-Unterstützung für das effiziente ADC-Sampling der Temperatursensoren nutzen. & Performance \\
		\addlinespace
		ID-0015 & Das System muss die Spannungswandlung von 12V auf 3,3V berücksichtigen und stabile Betriebsbedingungen für die MCU sicherstellen. & Design \\
		\addlinespace
		ID-0016 & Die Software muss eine Fehlerbehandlung für die ADC-Kanäle implementieren (z.B. bei Ausfall eines Temperatursensors). & Funktional \\
		\bottomrule
	\end{tabularx}
\end{table}


\subsection{Hardware- und Softwareanforderungen}
% Hier eine kurze Tabelle oder Liste der ca. 7 Requirements (z.B. REQ-1: Zwei Temperaturen messen, REQ-2: PWM-Steuerung für Pumpe).

\subsection{Reflexion zur Anforderungsanalyse}
% Ein kurzer Absatz, warum im Rahmen des Wasserfallmodells und der Projektgröße die Anforderungen pragmatisch gehalten wurden.


% ------------------------------------------
% 2. Systemarchitektur und Hardware-Design
% ------------------------------------------
\section{Systemarchitektur und Hardware-Design}
% Hier die Systemarchitektur beschreiben[cite: 197].

\subsection{Gesamtarchitektur}
% Platz für dein angepasstes PlantUML-Diagramm.
% \begin{figure}[h]
	%     \centering
	%     \includegraphics[width=0.8\textwidth]{bilder/architektur.png}
	%     \caption{Systemarchitektur des Bierkühlers}
	% \end{figure}

\subsection{Wahl des Mikrocontrollers und Spannungsversorgung}
% Warum der STM32? Erklärung der 3,3V / 12V Versorgungskomponenten.

\subsection{Sensorik (Eingabe)}
% Beschreibung der beiden Thermometer und deren Anbindung über den ADC.

\subsection{Aktorik (Ausgabe) und Technical Trades}
% Technical Trades[cite: 93].
% SEHR WICHTIG: Hier erklären, warum das Relais durch PWM-Module für Pumpe und Motor ersetzt wurde (z.B. stufenlose Regelung).

\subsection{User Interface}
% Anbindung und Ansteuerung des Displays über I²C.


% ------------------------------------------
% 3. Software-Design und Implementierung
% ------------------------------------------
\section{Software-Design und Implementierung}
% Hier die Softwarearchitektur beschreiben[cite: 201].

\subsection{Software-Architektur und Main-Loop}
% Flussdiagramm eurer Software. Beschreibung, dass es sich um ein "MCT optimiertes Design" handelt[cite: 180].

\subsection{Peripherie-Setup und Treiber}
\subsubsection{I²C-Kommunikation (Display)}
\subsubsection{PWM-Generierung (Motor und Pumpe)}
\subsubsection{ADC-Konfiguration (Temperatursensoren)}

\subsection{Code-Qualität und Dokumentation}
% Verweis auf "Applicable Standards and Rules (Coding Rules)"[cite: 206].
% Kurzes Beispiel zeigen, wie ihr Doxygen für die Code-Dokumentation genutzt habt[cite: 170, 258].


% ------------------------------------------
% 4. Test und Validierung
% ------------------------------------------
\section{Test und Validierung}
% Test Specification and Report[cite: 190].

\subsection{Testfälle (Test Cases)}
% Wie wurden die ~7 Requirements getestet? [cite: 195]

\subsection{Integrationstests der Hardware}
% Liefen Pumpe, Motor, Display und Sensoren fehlerfrei parallel am STM32?

\subsection{Testresultate}
% Ergebnisse der Tests dokumentieren[cite: 199].


% ------------------------------------------
% 5. Fazit und Lessons Learned
% ------------------------------------------
\section{Fazit und Lessons Learned}
% Dieser Teil ist extrem wichtig für die Präsentation und Dokumentation[cite: 176, 230].

\subsection{Erreichte Projektziele}
% Kurzes Fazit: Kühlt das System das Bier wie gewünscht?

\subsection{Lessons Learned: Was gut lief}
% Offene Kommunikation darüber, was im Team und technisch gut funktioniert hat[cite: 230].

\subsection{Lessons Learned: Herausforderungen und Fehler}
% Was hat weniger gut funktioniert und warum? [cite: 231] (z.B. Probleme beim Wasserfallmodell, späte Änderungen von Relais auf PWM).